\chapter*{\textbf{Zusammenfassung}}
\label{chap:abstract-german}
Die Bachelorarbeit befasst sich mit der Nutzung von generativer künstlicher Intelligenz in Bauprojekten, die datengetriebene Planungsmethoden einsetzen.
Dabei liegt der Fokus auf der Extraktion von Informationen aus im Bauprojekt vorhandenen Datenquellen durch Probanden mittels einer natürlichen Konversation. 
Die Probanden führen die Konversation mit einer im Rahmen der Arbeit entwickelten Textein- und -ausgabe, welche mit einem exemplarischen digitalen Modell verknüpft ist.
Die Probandengruppe umfasst neben Spezialisten, die einige Jahre Erfahrung mit der datengetriebenen Planung besitzen, gleichermaßen Personen ohne direkte Erfahrung im genannten Bereich. 
Bei der Analyse der künstlich generierten Inhalte wird die Richtigkeit der Ergebnisse im Kontext des verbundenen digitalen Modells beurteilt.
Abschließend werden die gewonnenen Ergebnisse eingeordnet und bewertet sowie ein Ausblick auf zukünftige Entwicklungen gegeben.


\chapter*{\textbf{Abstract}}
\label{chap:abstract-english}
\selectlanguage{english}
This bachelor thesis addresses the use of generative artificial intelligence in construction projects that utilize data-driven planning methods.
The primary focus lies on extracting information from data sources embedded within construction projects by subjects using a custom-developed text input and output interface during a natural conversation.
The subject group contains specialists with multiple years of experience in data-driven planning, as well as subjects without relevant prior experience.
When analyzing the artificially generated content, the accuracy of the generated responses is crucial in the context of the connected digital model.
Lastly, the results are categorized and evaluated, followed by an outlook on future developments.
\selectlanguage{german}